% GENERATED FILE. Edit canonical lessons, then run npm run generate:books.
% canonical-source-hash: fnv1a64:642bde13a1c606e9
% canonical-lessons: FR-C324-envahir, FR-C324-invasion, FR-C324-conquerir, FR-C324-conquete, FR-C324-vaincre

\chapter{La conquête}
\label{ch:fr-la-conquete}
\hlchaptermodality{324}

\begin{chapteropening}
\textbf{By the end of this chapter:} \emph{I can talk about wars of long ago: invading, an invasion, conquering, a conquest and defeating an enemy.}

Complete the last lesson of chapter 324: I can talk about wars of long ago: invading, an invasion, conquering, a conquest and defeating an enemy.
\end{chapteropening}

\section[envahir]{envahir — to invade; to overrun}

\label{lesson:FR-C324-envahir}



\begin{quote}
Before the new one: say the French for the crown, then the French for to crown.
\end{quote}

\subsection*{You'll want to know: envahir}
\textbf{envahir} — \textquotedblleft{}to invade; to overrun\textquotedblright{}.

\begin{quote}
L'armée a envahi le pays voisin.
\end{quote}

\begin{quote}
\emph{The army invaded the neighbouring country.}
\end{quote}

Figurative: \emph{les touristes envahissent la plage}.

\begin{cousinweb}
\emph{envahir} is inherited from Vulgar Latin \emph{invādīre}, from Latin \emph{invādere} \textquotedblleft{}to go into, to attack\textquotedblright{} (EtymDB).
\end{cousinweb}

\begin{grammarlens}[title={What you've built}]
One more word for telling, explaining and describing at B1.
\end{grammarlens}

\subsection*{Your turn}
\begin{itemize}
\raggedright
  \item \emph{Say it:} \emph{envahir}
  \item \emph{Say it:} \emph{envahir}, once more
  \item \emph{Say it:} a sentence of your own with \emph{envahir}
  \item \emph{Recall:} say the French for the crown, then the French for to crown, then say \emph{envahir} again
\end{itemize}

\begin{tcolorbox}[breakable,colback=teal!4,colframe=teal!35!black,title={Before you move on}]
What does envahir mean? (\textquotedblleft{}To invade; to overrun\textquotedblright{}.) Say it once more.
\end{tcolorbox}
\section[l'invasion]{l'invasion — the invasion}

\label{lesson:FR-C324-invasion}



\begin{quote}
Before the new one: say the French for to crown, then the French for to invade.
\end{quote}

\subsection*{You'll want to know: l'invasion}
\textbf{l'invasion} — \textquotedblleft{}the invasion\textquotedblright{}.

\begin{quote}
L'invasion a commencé au printemps.
\end{quote}

\begin{quote}
\emph{The invasion began in the spring.}
\end{quote}

Figurative: \emph{une invasion de moustiques}.

\begin{cousinweb}
\emph{invasion} is a learned borrowing from Late Latin \emph{invāsiō} (EtymDB), from \emph{invādere} \textquotedblleft{}to go into\textquotedblright{}.
\end{cousinweb}

\begin{grammarlens}[title={What you've built}]
One more word for telling, explaining and describing at B1.
\end{grammarlens}

\subsection*{Your turn}
\begin{itemize}
\raggedright
  \item \emph{Say it:} \emph{l'invasion}
  \item \emph{Say it:} \emph{l'invasion}, once more
  \item \emph{Say it:} a sentence of your own with \emph{l'invasion}
  \item \emph{Recall:} say the French for to crown, then the French for to invade, then say \emph{l'invasion} again
\end{itemize}

\begin{tcolorbox}[breakable,colback=teal!4,colframe=teal!35!black,title={Before you move on}]
What does l'invasion mean? (\textquotedblleft{}The invasion\textquotedblright{}.) Say it once more.
\end{tcolorbox}
\section[conquérir]{conquérir — to conquer; to win over}

\label{lesson:FR-C324-conquerir}



\begin{quote}
Before the new one: say the French for to invade, then the French for the invasion.
\end{quote}

\subsection*{You'll want to know: conquérir}
\textbf{conquérir} — \textquotedblleft{}to conquer; to win over\textquotedblright{}.

\begin{quote}
Ils ont conquis toute la région en un an.
\end{quote}

\begin{quote}
\emph{They conquered the whole region in a year.}
\end{quote}

Past participle \emph{conquis}. Also \emph{conquérir le public} — to win the audience over.

\begin{cousinweb}
\emph{conquérir} comes from Old French \emph{conquerre}, from Vulgar Latin \emph{conquaerere} \textquotedblleft{}to seek to gain\textquotedblright{} (EtymDB), built on \emph{quaerere} \textquotedblleft{}to seek\textquotedblright{}.
\end{cousinweb}

\begin{grammarlens}[title={What you've built}]
One more word for telling, explaining and describing at B1.
\end{grammarlens}

\subsection*{Your turn}
\begin{itemize}
\raggedright
  \item \emph{Say it:} \emph{conquérir}
  \item \emph{Say it:} \emph{conquérir}, once more
  \item \emph{Say it:} a sentence of your own with \emph{conquérir}
  \item \emph{Recall:} say the French for to invade, then the French for the invasion, then say \emph{conquérir} again
\end{itemize}

\begin{tcolorbox}[breakable,colback=teal!4,colframe=teal!35!black,title={Before you move on}]
What does conquérir mean? (\textquotedblleft{}To conquer; to win over\textquotedblright{}.) Say it once more.
\end{tcolorbox}
\section[la conquête]{la conquête — the conquest}

\label{lesson:FR-C324-conquete}



\begin{quote}
Before the new one: say the French for the invasion, then the French for to conquer.
\end{quote}

\subsection*{You'll want to know: la conquête}
\textbf{la conquête} — \textquotedblleft{}the conquest\textquotedblright{}.

\begin{quote}
La conquête de l'espace a commencé en 1957.
\end{quote}

\begin{quote}
\emph{The conquest of space began in 1957.}
\end{quote}

Informal: \emph{sa dernière conquête} — his latest girlfriend or boyfriend.

\begin{cousinweb}
\emph{conquête} comes from Old French \emph{conqueste}, usually traced to Vulgar Latin \emph{conquaesita} \textquotedblleft{}things sought and won\textquotedblright{}, built on \emph{quaerere} \textquotedblleft{}to seek\textquotedblright{} (Lewis \& Short).
\end{cousinweb}

\begin{grammarlens}[title={What you've built}]
One more word for telling, explaining and describing at B1.
\end{grammarlens}

\subsection*{Your turn}
\begin{itemize}
\raggedright
  \item \emph{Say it:} \emph{la conquête}
  \item \emph{Say it:} \emph{la conquête}, once more
  \item \emph{Say it:} a sentence of your own with \emph{la conquête}
  \item \emph{Recall:} say the French for the invasion, then the French for to conquer, then say \emph{la conquête} again
\end{itemize}

\begin{tcolorbox}[breakable,colback=teal!4,colframe=teal!35!black,title={Before you move on}]
What does la conquête mean? (\textquotedblleft{}The conquest\textquotedblright{}.) Say it once more.
\end{tcolorbox}
\section[vaincre]{vaincre — to defeat, to overcome}

\label{lesson:FR-C324-vaincre}



\begin{quote}
Before the new one: say the French for to conquer, then the French for the conquest.
\end{quote}

\subsection*{You'll want to know: vaincre}
\textbf{vaincre} — \textquotedblleft{}to defeat, to overcome\textquotedblright{}.

\begin{quote}
L'armée a vaincu l'ennemi en trois jours.
\end{quote}

\begin{quote}
\emph{The army defeated the enemy in three days.}
\end{quote}

Irregular: \emph{il a vaincu}, \emph{nous vainquons}.

\begin{cousinweb}
\emph{vaincre} comes from Old French \emph{veintre} (EtymDB), inherited from Latin \emph{vincere} \textquotedblleft{}to conquer\textquotedblright{} (Lewis \& Short).
\end{cousinweb}

\begin{grammarlens}[title={What you've built}]
One more word for telling, explaining and describing at B1.
\end{grammarlens}

\subsection*{Your turn}
\begin{itemize}
\raggedright
  \item \emph{Say it:} \emph{vaincre}
  \item \emph{Say it:} \emph{vaincre}, once more
  \item \emph{Say it:} a sentence of your own with \emph{vaincre}
  \item \emph{Recall:} say the French for to conquer, then the French for the conquest, then say \emph{vaincre} again
\end{itemize}

\begin{tcolorbox}[breakable,colback=teal!4,colframe=teal!35!black,title={Before you move on}]
What does vaincre mean? (\textquotedblleft{}To defeat, to overcome\textquotedblright{}.) Say it once more.
\end{tcolorbox}
